\documentclass[11pt,a4paper]{article}

%% =====================================================================
%% Clay P vs NP via Substrate --- Phase 3 (v2, HONEST-SCOPE REWRITE)
%%
%% Title: Phase 3 --- A Substrate-Level Reformulation at the
%%        P vs NP Axis via the Principia Fractalis Framework
%%
%% Author: Pablo Cohen (with Claude Opus 4.7 / Anthropic)
%% Date  : 2026-06-06
%% Anchor: HEAD 4785c55 / cc7d9d3, Lean 4 PF closure clean,
%%         zero project axioms, kernel-only
%%         [propext, Classical.choice, Quot.sound].
%% Honest-scope discipline:
%%   Every claim sourceable to a Lean file's own honest-scope.
%%   The framework's `Machine` type has free functions and no
%%   operational Turing semantics; this is foregrounded.
%% =====================================================================

\usepackage{amsmath,amsthm,amssymb,amsfonts}
\usepackage[margin=1in]{geometry}
\usepackage{hyperref}
\usepackage{cleveref}
\usepackage{microtype}
\usepackage{enumitem}
\usepackage{booktabs}
\usepackage{xcolor}

\hypersetup{
  colorlinks=true,
  linkcolor=blue!50!black,
  urlcolor=blue!50!black,
  citecolor=blue!50!black,
}

\theoremstyle{plain}
\newtheorem{theorem}{Theorem}[section]
\newtheorem{lemma}[theorem]{Lemma}
\newtheorem{proposition}[theorem]{Proposition}
\newtheorem{corollary}[theorem]{Corollary}

\theoremstyle{definition}
\newtheorem{definition}[theorem]{Definition}
\newtheorem{solution}[theorem]{Solution}
\newtheorem{remark}[theorem]{Remark}

\title{Phase 3: A Substrate-Level Reformulation at the
P vs NP Axis\\
via the Principia Fractalis Framework\\
\large (Not a Clay Discharge; Machine Type Carries No Operational
Turing Semantics)}
\author{Pablo Cohen\\
\small{\texttt{psolorzano@gmail.com}}\\[0.3em]
\small{Collaboration: Claude Opus 4.7 (Anthropic)}}
\date{2026-06-06}

\begin{document}
\maketitle

\begin{abstract}
This is the Phase 3 paper of a four-paper sequence on the
Principia Fractalis (PF) framework's substrate-level reformulation
of the six remaining Clay Millennium Problems. \textbf{This is not a
Clay-prize claim for P vs NP.} The framework's content on this axis
is: (i)~the chosen value $\alpha_{\mathsf{P}\mathrm{vs}\mathsf{NP}}
:= 5/4$ pinned by the eleven framework-internal arithmetic
compatibility identities on the chosen $\alpha$-skeleton under the
editorial convention $\alpha_{\mathrm{Poincar\acute e}} := 1$
motivated by Perelman 2003; (ii)~an unconditional biconditional
$\mathtt{Clay\_PvsNP\_Standard\;PF\_CanonicalComplexityEncoding}
\leftrightarrow (\mathtt{ClassP} \neq \mathtt{ClassNP})$ on the
framework's chosen complexity-encoding; (iii)~the open Clay
conjecture itself ($\mathtt{ClassP} \neq \mathtt{ClassNP}$ on the
chosen encoding) as the named open residual the framework does
\emph{not} discharge.

The framework's $\mathtt{Machine}$ type is \emph{not} an operational
Turing machine. Its two fields ($\mathtt{accepts : BinString \to
Prop}$ and $\mathtt{steps : BinString \to \mathbb{N}}$) are free
functions; the three type parameters $\Gamma, \Lambda, \sigma$ are
unused in the body. The file docstring acknowledges this; the
classes $\mathtt{ClassP}$ and $\mathtt{ClassNP}$ built atop this
structure are substrate-shadows of the complexity-theoretic
structure, not the operational Cook--Karp form. The framework's
contribution on this axis is the biconditional and the
$\alpha$-skeleton pinning; it is \emph{not} a constructive
polynomial-time algorithm for SAT, and it is \emph{not} a literal
mathlib-tier discharge of $\mathsf{P} \neq \mathsf{NP}$.

Machine-verified in Lean 4 with kernel-only axioms
$\mathtt{[propext, Classical.choice, Quot.sound]}$. The
PF\_Lean4Lean harness provides build-hash robustness across two
Lake packages (\emph{not} an independent type-checker in another
language). Honest scope is preserved throughout. Independent expert
review and substantial revision are explicitly invited.
\end{abstract}

\tableofcontents

\section{What this paper offers and what it does not}
\label{sec:offers}

This paper offers a \emph{substrate-level reformulation of the
P vs NP axis} of the Principia Fractalis framework. It does
\emph{not} offer a literal mathlib-tier discharge of $\mathsf{P} \neq
\mathsf{NP}$ in the Clay form, and it does \emph{not} offer a
constructive polynomial-time algorithm for SAT or for any other
NP-complete problem. The substrate-level content is the unconditional
biconditional on the framework's chosen encoding and the
$\alpha$-skeleton pinning of $\alpha_{\mathsf{P}\mathrm{vs}\mathsf{NP}}
:= 5/4$, $\alpha_{P} := \sqrt{2}$, $\alpha_{NP} := \varphi + 1/4$;
the open Clay conjecture $\mathtt{ClassP} \neq \mathtt{ClassNP}$ on
the chosen encoding is preserved openly.

\medskip

\noindent\fbox{\parbox{0.97\textwidth}{
\textbf{What the Lean kernel carries on this axis.} The biconditional
$\mathtt{Clay\_PvsNP\_Standard\;PF\_CanonicalComplexityEncoding}
\leftrightarrow (\mathtt{ClassP} \neq \mathtt{ClassNP})$ unconditional
on the framework's chosen encoding; the chosen
$\alpha_{\mathsf{P}\mathrm{vs}\mathsf{NP}} := 5/4$ value as the unique
solution to the framework's algebraic system under the editorial
convention; the bundle field
$\mathtt{pnp\_classes\_distinct} : \mathtt{ClassP} \neq
\mathtt{ClassNP}$ from the simultaneous-closure bundle of Phase 2
(this field \emph{is} the P vs NP question on the chosen encoding,
preserved openly).

\textbf{What this paper does NOT claim.} A literal mathlib-tier
discharge of $\mathsf{P} \neq \mathsf{NP}$ (the form required by the
Clay rules). A constructive polynomial-time algorithm for any
NP-complete problem. A separation through a small-polynomial
separator. The framework's $\mathtt{Machine}$ type carries no
operational Turing semantics; on this encoding, the discharge is a
substrate-shadow of the Cook--Karp operational form, not the
operational form itself.
}}

\section{The framework's $\mathtt{Machine}$ type, foregrounded}
\label{sec:machine-honest}

The framework's complexity-axis encoding lives in
$\mathtt{PrincipiaTractalis.TuringEncoding.Complexity}$. Its
$\mathtt{Machine}$ structure is defined (line~71) as:

\begin{verbatim}
structure Machine (Gamma Lambda sigma : Type) where
  accepts : BinString -> Prop
  steps   : BinString -> N
\end{verbatim}

\noindent
\textbf{Honest scope, foregrounded:}
\begin{itemize}[leftmargin=*,nosep]
\item The three type parameters $\Gamma, \Lambda, \sigma$ are
  \emph{unused} in the body of the structure.
\item The two fields $\mathtt{accepts}$ and $\mathtt{steps}$ are
  \emph{free functions} with no operational semantics constraining
  them to be Turing-computable.
\end{itemize}
The file docstring at the structure (lines 57--70 of
\texttt{Complexity.lean}) acknowledges this: ``\emph{the full mathlib}
$\mathtt{Turing.TM2}$ \emph{model parametrizes over} $K$ \emph{(stack
indices)}, $\Gamma : K \to \mathtt{Type}$ \emph{(alphabet per stack)},
$\Lambda$ \emph{(state set), and} $\sigma$ \emph{(auxiliary register).
Here we keep the surface signature\ldots}''

The classes $\mathtt{ClassP}$ (line~118) and $\mathtt{ClassNP}$
(line~158) built atop this structure are therefore
\emph{substrate-shadows} of the Cook 1971 / Karp 1972
operational form, not the operational form itself. The framework's
$\mathtt{InClassP}$ predicate says ``there exists a machine $M$ with
$\mathtt{steps}$ polynomially bounded and $\mathtt{accepts}$ agreeing
with membership in $L$''; but the existential ranges over the
free-function $\mathtt{Machine}$ population, not over operational
Turing machines, so its meaning is the framework's substrate-shadow
of the operational class.

\textbf{What this means for the framework's P vs NP claim.} The
biconditional and the $\mathtt{ClassP} \neq \mathtt{ClassNP}$
formulation are correctly stated \emph{on the framework's chosen
encoding}. They are not the literal Cook--Karp operational
formulation, and the framework does not claim they are. A literal
mathlib-tier discharge would require either porting the framework's
encoding to mathlib's $\mathtt{Turing.TM2}$ or constructing a
specific machine inhabiting the predicates; neither is done here. The
open Clay conjecture, in the framework's chosen encoding, is the
bundle field $\mathtt{pnp\_classes\_distinct}$ of Phase 2's
seven-field bundle.

\section{The substrate and the editorial convention}
\label{sec:substrate}

The framework's substrate is the Timeless Field
$\mathcal{H}_{\mathrm{TF}} = \varprojlim_{k} \mathbb{C}^{3^{k}}$
(see the master paper \cite{master_substrate} for the substrate
description). The editorial convention
$\alpha_{\mathrm{Poincar\acute e}} := 1$ is motivated by Perelman
2003's resolution of the Poincar\'e conjecture \cite{perelman2002,%
perelman2003a,perelman2003b}; the Lean kernel does not consume
Perelman's theorem as a hypothesis. The convention is the framework's
identification of the substrate's Poincar\'e-axis constant.

The eleven framework-internal arithmetic compatibility identities on
the chosen $\alpha$-skeleton are catalogued in the master paper and
in Phase 2. Honest scope, quoted verbatim from
\texttt{CrossMillenniumSharedInvariants.lean} (lines 18--29):
``\emph{These are \textbf{not} Millennium discharges. They are
axiom-free algebraic facts about the 9 chosen $\alpha$-values; the
framework's interpretation\ldots is conjectural and lives elsewhere.}''

\section{The $\alpha$-values on the complexity axis}
\label{sec:alpha-complexity}

The framework's $\alpha$-skeleton has three distinguished values on
the complexity axis: $\alpha_{P} := \sqrt{2}$, $\alpha_{NP} := \varphi
+ 1/4$ (where $\varphi = (1 + \sqrt{5})/2$), and
$\alpha_{\mathsf{P}\mathrm{vs}\mathsf{NP}} := 5/4$, each defined in
$\mathtt{PF.CrossMillenniumSharedInvariants}$.

\begin{theorem}[Chosen complexity-axis $\alpha$-values pinned by the
framework-internal identities under the editorial convention]
\label{thm:forced-complexity}
The chosen complexity-axis values
\[
\alpha_{P} = \sqrt{2}, \qquad
\alpha_{NP} = \varphi + 1/4, \qquad
\alpha_{\mathsf{P}\mathrm{vs}\mathsf{NP}} = 5/4
\]
are the unique solution to the framework's algebraic system under the
editorial convention $\alpha_{\mathrm{Poincar\acute e}} := 1$ and the
eleven framework-internal arithmetic identities (e.g.\
$\alpha_{P}^{2} = \alpha_{\mathrm{YM}}$ where
$\alpha_{\mathrm{YM}} = \alpha_{\mathrm{Poincar\acute e}} + 1 = 2$;
$\alpha_{NP} - \alpha_{\mathrm{Hodge}} = 1/4$ where
$\alpha_{\mathrm{Hodge}} = \varphi$).
\textbf{Lean:}
\texttt{PF.Referee.ClayMasterTheorem.\\
framework\_alpha\_unique\_under\_perelman\_anchor};
\texttt{PrincipiaTractalis.CrossMillenniumSharedInvariants.\\
cross\_millennium\_shared\_invariants\_capstone};
$\alpha\_P\_sq\_eq\_\alpha\_YM$;
$\alpha\_NP\_sub\_Hodge\_eq\_quarter$;
$\alpha\_Hodge\_sq\_eq\_self\_plus\_one$;
$\alpha\_YM\_eq\_\alpha\_Poincare\_plus\_one$.
\end{theorem}

\noindent\textbf{Honest scope.} These are arithmetic facts about
\emph{chosen} constants under the editorial convention. The
framework's interpretation that $\alpha_{P}$ and $\alpha_{NP}$ are
the canonical eigenvalue parameters of the $\mathsf{P}$ and
$\mathsf{NP}$ operators in the framework's polylog encoding is, per
\texttt{CrossMillenniumSharedInvariants.lean} header, ``conjectural
and lives elsewhere.''

\section{The canonical complexity encoding and the biconditional}
\label{sec:canonical-encoding}

\subsection{The framework's chosen complexity encoding}
\label{sec:chosen-encoding}

The framework's referee module $\mathtt{StandardClayStatements}$
parameterises the Clay P vs NP statement over an external encoding
$\mathtt{StandardComplexityEncoding}$ (three fields:
$\mathtt{ClassP : Type}$, $\mathtt{ClassNP : Type}$,
$\mathtt{inclusion : ClassP \to ClassNP}$). The Clay-statement
contract on an encoding $E$ is
\[
\mathtt{Clay\_PvsNP\_Standard}\, E := \neg\,
\mathtt{Function.Surjective}\, E.\mathtt{inclusion}.
\]
The framework's chosen instantiation is:

\begin{definition}[The framework's chosen complexity-encoding]
\label{def:pf-canonical-encoding}
\begin{align*}
\mathtt{CanonicalClassP} &:= \{ L : \mathtt{Language} \mid L \in
\mathtt{ClassP} \}, \\
\mathtt{CanonicalClassNP} &:= \{ L : \mathtt{Language} \mid L \in
\mathtt{ClassNP} \}, \\
\mathtt{PF\_CanonicalComplexityEncoding} &:= \langle
\mathtt{CanonicalClassP}, \mathtt{CanonicalClassNP},
\mathtt{canonical\_PvsNP\_inclusion} \rangle.
\end{align*}
The inclusion $\mathtt{canonical\_PvsNP\_inclusion} :
\mathtt{CanonicalClassP} \to \mathtt{CanonicalClassNP}$ is defined by
$L \mapsto \langle L.\mathtt{val}, \mathtt{P\_subset\_NP}\,
L.\mathtt{property}\rangle$, where
$\mathtt{P\_subset\_NP}$ is the framework's typed proof that the
framework's $\mathtt{ClassP}$ is included in
$\mathtt{ClassNP}$. \textbf{Lean:}
\texttt{PF.Referee.PNPCanonicalEncoding.\\
PF\_CanonicalComplexityEncoding} (line~81).
\end{definition}

\noindent\textbf{Honest scope, foregrounded again.} This encoding is
built from the framework's $\mathtt{ClassP}$/$\mathtt{ClassNP}$,
which are in turn built from the framework's $\mathtt{Machine}$ type
whose fields are free functions with no operational Turing semantics
(\S\ref{sec:machine-honest}). The encoding is \emph{the framework's
substrate-shadow} of the Cook--Karp form. Earlier framework documents
that called it ``the literal mathlib-canonical Cook 1971 / Karp 1972
encoding'' overstated the encoding's status; it is the framework's
chosen encoding built atop a $\mathtt{Machine}$ structure that does
not carry operational Turing semantics. The framework's commitment to
the Cook--Karp form is structural (subtype quotient of languages
satisfying the predicates) but is not validated through an
operational-machine layer.

\subsection{The unconditional biconditional}
\label{sec:biconditional}

\begin{theorem}[Biconditional on the chosen encoding]
\label{thm:canonical-iff-distinct}
On the framework's chosen complexity-encoding,
\[
\mathtt{Clay\_PvsNP\_Standard}\,
\mathtt{PF\_CanonicalComplexityEncoding}
\;\Longleftrightarrow\;
\mathtt{ClassP} \neq \mathtt{ClassNP}.
\]
\textbf{Lean:}
\texttt{PF.Referee.PNPCanonicalEncoding.\\
Clay\_PvsNP\_Standard\_at\_canonical\_iff\_classes\_distinct}
(line~92). Kernel axioms:
\texttt{[propext, Classical.choice, Quot.sound]}.
\end{theorem}

\noindent The biconditional itself is unconditional on the chosen
encoding. \textbf{The substrate-level content is the biconditional;
the right-hand side ($\mathtt{ClassP} \neq \mathtt{ClassNP}$ on the
chosen encoding) is the open Clay conjecture, preserved openly}, and
appears as the bundle field $\mathtt{pnp\_classes\_distinct}$ in the
Phase 2 simultaneous-closure bundle.

\section{The polylog typed upgrade --- conditional content}
\label{sec:polylog-typed}

The framework's $\mathtt{PolylogEigenvalueConjecture}$ residual
(\texttt{PF.TuringEncoding.Operators.PolylogEigenvalueConjecture})
is the existential statement that there exists a class function
$f : \mathtt{Set\;Language} \to \mathbb{R}$ satisfying the
framework's four typed sub-Props (P-class self-adjointness, P-class
positivity, NP-class golden-modulation quadratic, NP-class
positivity). The framework's typed biconditional reduces this
existential to $\mathtt{ClassP} \neq \mathtt{ClassNP}$:

\begin{theorem}[Polylog typed upgrade: existential $\Leftrightarrow$
class distinctness]
\label{thm:polylog-typed-upgrade}
\[
\bigl(\exists f, \mathrm{A1}(f) \wedge \mathrm{A2}(f) \wedge
\mathrm{B1}(f) \wedge \mathrm{B2}(f)\bigr)
\;\Longleftrightarrow\;
\mathtt{ClassP} \neq \mathtt{ClassNP}.
\]
\textbf{Lean:}
\texttt{PrincipiaTractalis.TuringEncoding.\\
polylog\_subprop\_quadruple\_iff\_classes\_distinct}.
\end{theorem}

\noindent\textbf{Honest scope.} The biconditional is unconditional;
its content is logical equivalence of the existential and class
distinctness on the framework's chosen encoding, \emph{not} a proof
of either. The four sub-Props are discharged axiom-free \emph{at the
substrate enum-level} (the framework's enumeration on the chosen
$(\alpha_{P}, \alpha_{NP}) = (\sqrt{2}, \varphi + 1/4)$), per
$\mathtt{polylog\_subprop\_quadruple\_at\_enum}$; this is a
substrate-level discharge, not a literal mathlib-tier discharge.

\section{Connection to the simultaneous-closure mechanism}
\label{sec:simultaneous-pointer}

The substrate-level Clay-Standard contract on the chosen P vs NP
encoding is one of six axes in the Phase 2 simultaneous-closure
mechanism. The bundle field $\mathtt{pnp\_classes\_distinct} :
\mathtt{ClassP} \neq \mathtt{ClassNP}$ is the typed residual
consumed for this axis; the simultaneous-closure theorem
$\mathtt{perelman\_anchor\_yields\_simultaneous\_clay\_closure}$
delivers the six Clay-Standards on canonical encodings under the
editorial convention plus the seven-field bundle (5 typed residuals
+ 2 $\mathtt{True}$-markers).

The framework's contribution at this axis is structural: the
biconditional and the $\alpha$-skeleton pinning on the chosen
encoding; the open Clay conjecture itself is preserved openly. See
the master paper \cite{master_substrate} for the simultaneous-closure
mechanism's full statement and per-axis honest-scope discussion.

\section{Independent verification posture}
\label{sec:l4l}

The Lean~4 proof object for the theorems cited type-checks in the
Lean~4 kernel with kernel axioms $\mathtt{[propext, Classical.choice,
Quot.sound]}$ and zero project axioms. The PF\_Lean4Lean harness
provides \textbf{build-hash robustness across two Lake packages}:
the canonical theorems are re-bound in a separate package and
rebuilt from source. \textbf{It is not an independent type-checker
written in another language.} Quoted verbatim from
\texttt{PF\_Lean4Lean/PF\_L4L/Referee/ClayVerificationHarness.lean}
lines~47--56:
\begin{quote}
``\emph{This is the same Path-C protocol\ldots It is NOT a separate
type-checker written in another language. The honest scope of each
individual closure is the one fixed by that closure's own
\texttt{*\_honestScope} / \texttt{*\_honest\_scope} marker in the
canonical PF library\ldots The harness re-exports the closures as
they are --- it does NOT strengthen any of them.}''
\end{quote}

\textbf{Earlier drafts called this an ``independent kernel
re-verifier''; that framing overstated the harness's content and is
dropped in favor of the file's own self-description.}

\section{What this does and does not open}
\label{sec:what-opens}

The framework's contribution on the P vs NP axis is structural at the
$\alpha$-skeleton level on the chosen encoding. \textbf{The
framework does not supply:}
\begin{itemize}[leftmargin=*,nosep]
\item A constructive polynomial-time algorithm for SAT or for any
  other NP-complete problem.
\item A literal mathlib-tier discharge of $\mathsf{P} \neq
  \mathsf{NP}$.
\item A small-polynomial separator.
\item A bypass of Razborov--Rudich 1997 natural-proofs
  \cite{razborov_rudich1997} or Aaronson--Wigderson 2009
  algebrisation \cite{aaronson_wigderson2009} barriers.
\end{itemize}
What the framework supplies is the substrate-level reformulation: the
$\alpha$-skeleton pinning of $\alpha_{\mathsf{P}\mathrm{vs}\mathsf{NP}}
= 5/4$ on the chosen encoding, the biconditional reducing the
Clay-Standard contract to the open class-distinctness conjecture on
the chosen encoding, and the integration of this axis into the
simultaneous-closure mechanism of Phase 2 alongside the other five
remaining Clay axes.

The deployable infrastructure that operates today on the working
assumption $\mathsf{P} \neq \mathsf{NP}$ (cryptographic primitives,
TLS handshakes, SAT solvers as heuristics, integer-programming
packages, logistics optimizers) is not affected by the framework's
substrate-level content. The framework neither validates nor
invalidates that assumption at the operational tier.

\section{Honest scope}
\label{sec:honest-scope}

Per the file-level honest-scope marker
($\mathtt{sec:substrate-honest-scope}$), PF is a
\emph{substrate-level monograph in Grothendieck mode}.

\begin{remark}[Honest scope at the P vs NP axis, foregrounded]
\label{rem:honest-scope-pnp}

\textbf{(H1) Not a literal mathlib-tier Clay discharge.} The framework
does not discharge $\mathsf{P} \neq \mathsf{NP}$ at the form the Clay
rules require. The framework's biconditional is unconditional on the
chosen encoding; the right-hand side ($\mathtt{ClassP} \neq
\mathtt{ClassNP}$ on the chosen encoding) is the open Clay
conjecture, preserved openly.

\textbf{(H2) The $\mathtt{Machine}$ type carries no operational Turing
semantics.} The framework's $\mathtt{Machine}$ structure (line~71 of
\texttt{Complexity.lean}) has two fields ($\mathtt{accepts}$,
$\mathtt{steps}$) that are free functions; three type parameters
($\Gamma, \Lambda, \sigma$) unused in the body. The file docstring
acknowledges this. The classes $\mathtt{ClassP}$ and $\mathtt{ClassNP}$
built atop this structure are substrate-shadows of the Cook--Karp
operational form, not the operational form itself.

\textbf{(H3) Not a constructive polynomial-time SAT solver.} The
framework does not supply a polynomial-time algorithm for any
NP-complete problem. The framework makes no claim of constructive
algorithmic content.

\textbf{(H4) $\alpha$-skeleton pinning is framework-internal.} The
chosen values $\alpha_{P} := \sqrt{2}$, $\alpha_{NP} := \varphi +
1/4$, $\alpha_{\mathsf{P}\mathrm{vs}\mathsf{NP}} := 5/4$ are pinned
by the framework-internal arithmetic identities under the editorial
convention. The interpretation that these are the canonical
eigenvalue parameters of the $\mathsf{P}$ and $\mathsf{NP}$
operators is, per \texttt{CrossMillenniumSharedInvariants.lean}
header, ``conjectural and lives elsewhere''; the algebraic forcing of
the chosen values is internal to the framework's algebraic system,
not the consumption of external complexity-theoretic content.

\textbf{(H5) Cross-Millennium NP $\leftrightarrow$ Hodge bridge is an
identity on chosen constants.} The framework's identity
$\alpha_{NP} - \alpha_{\mathrm{Hodge}} = 1/4$ is an algebraic identity
on the framework's chosen $\alpha$-values; it is not a derivation of
either constant from external content.

\textbf{(H6) The PF\_Lean4Lean harness is build-hash robustness.} As
\S\ref{sec:l4l} states.
\end{remark}

\begin{remark}[The 143-problem coherence figure is retracted as
established evidence]
\label{rem:143-retraction}
The $\mathtt{universal\_fractal\_coherence}$ Lean theorem
(\texttt{PF.Empirical.HundredFortyThreeProblems}, line~167) records
the framework's prediction that all problems within a complexity
class P (resp.\ NP) yield identical canonical $\alpha$-values; the
dataset is implemented as
\texttt{List.replicate 72 (canonicalEntry .P) ++}
\texttt{List.replicate 71 (canonicalEntry .NP)}, and the
$\mathtt{coherence\_p\_value\_threshold}$ is asserted, not derived.
The file's own honest-scope note (lines 36--39) labels this as ``a
modeling assumption about a baseline success rate, not a derived
theorem.'' The $p < 10^{-43}$ figure should not be cited as
established empirical evidence; earlier framework documents that did
so are corrected here.
\end{remark}

\begin{remark}[IBM Quantum 9-way: pre-registration, 2 of 9 measured]
\label{rem:ibm-pre-registration}
At HEAD \texttt{4785c55} (\texttt{IBMHardware9WayEvidence.lean} lines
312--327), \textbf{only the pair $(\alpha_{\mathrm{RH}},
\alpha_{NP})$ has been observed on IBM Quantum hardware} consistent
with the framework's prediction. The other 7 are framework
predictions committed BEFORE measurement. The Lean theorem records a
\textbf{conditional} bound: ``IF all 9 are eventually measured to
within $10^{-3}$ of the pre-specified values under
$\mathtt{NoiseModel9}$, THEN random-match probability is at most
$10^{-15}$.'' Current status: 2 of 9 measured-consistent. The
pre-registration of seven framework $\alpha$-values prior to hardware
measurement is scientifically valuable; the conditional bound says
nothing absent the remaining 7 measurements.
\end{remark}

\begin{remark}[Falsifiability]
\label{rem:falsifiability}
The framework commits to eight typed falsifiability conditions in
\texttt{PF.Referee.FrameworkFalsifiabilityConditions}. Audit
2026-06-04: \textbf{zero of eight triggered}; \textbf{two actively
supported by 2024--2026 literature} ($F_{3}$, $F_{6}$).
\end{remark}

\section{Conclusion}
\label{sec:conclusion}

This Phase 3 paper presents the framework's substrate-level
reformulation at the P vs NP axis. The framework's contribution is
structural: an unconditional biconditional on the framework's chosen
encoding reducing $\mathtt{Clay\_PvsNP\_Standard}$ to
$\mathtt{ClassP} \neq \mathtt{ClassNP}$ on that encoding, an
$\alpha$-skeleton pinning of $\alpha_{\mathsf{P}\mathrm{vs}\mathsf{NP}}
:= 5/4$ under the editorial convention, and integration into the
Phase 2 simultaneous-closure mechanism. The framework's
$\mathtt{Machine}$ type carries no operational Turing semantics; the
chosen encoding is the framework's substrate-shadow of the Cook--Karp
form, not the operational form. \textbf{No Clay-prize claim is made.
No constructive polynomial-time algorithm is offered. The literal
$\mathsf{P} \neq \mathsf{NP}$ on the chosen encoding is the open
Clay conjecture, preserved openly.} Independent expert review and
substantial revision are explicitly invited.

\begin{thebibliography}{99}

\bibitem{master_substrate}
Cohen, P., Claude Opus 4.7 (2026-06-06). \emph{A Substrate-Level
Reformulation of the Six Remaining Clay Millennium Problems via the
Principia Fractalis Framework}.
\texttt{Papers/clay\_substrate\_reformulation\_v1.tex}.

\bibitem{phase1_rh}
Cohen, P., Claude Opus 4.7 (2026-06-06). \emph{Phase 1 ---
A Substrate-Level Reformulation of the Riemann Hypothesis via the
Principia Fractalis Framework}.
\texttt{Papers/clay\_RH\_via\_substrate\_v2.tex}.

\bibitem{phase2_six}
Cohen, P., Claude Opus 4.7 (2026-06-06). \emph{Phase 2 ---
The Simultaneous-Closure Mechanism on Canonical Encodings}.
\texttt{Papers/clay\_six\_simultaneous\_closure\_v1.tex}.

\bibitem{phase4_keystone}
Cohen, P., Claude Opus 4.7 (2026-06-06). \emph{Phase 4 Keystone ---
Substrate-Level Content on Consciousness, Zero-Point Energy, and
Cosmology}.
\texttt{Papers/principia\_fractalis\_keystone\_v1.tex}.

\bibitem{principia_book}
Cohen, P. (2026). \emph{Principia Fractalis}, Second Edition.
HEAD of
\url{https://github.com/FractalDevTeam/Principia-Fractalis}.

\bibitem{cook1971}
Cook, S.\ A. (1971). The complexity of theorem-proving procedures.
\emph{Proceedings of the Third Annual ACM Symposium on Theory of
Computing}, 151--158.

\bibitem{karp1972}
Karp, R. M. (1972). Reducibility among combinatorial problems.
In \emph{Complexity of Computer Computations}, Plenum, 85--103.

\bibitem{razborov_rudich1997}
Razborov, A. A., Rudich, S. (1997). Natural proofs.
\emph{J.\ Comput.\ System Sci.}\ 55(1): 24--35.

\bibitem{aaronson_wigderson2009}
Aaronson, S., Wigderson, A. (2009). Algebrization: a new barrier in
complexity theory. \emph{ACM Trans.\ Comput.\ Theory} 1(1), Art.\ 2.

\bibitem{perelman2002}
Perelman, G. (2002). The entropy formula for the Ricci flow and its
geometric applications. arXiv:math/0211159.

\bibitem{perelman2003a}
Perelman, G. (2003). Ricci flow with surgery on three-manifolds.
arXiv:math/0303109.

\bibitem{perelman2003b}
Perelman, G. (2003). Finite extinction time for the solutions to the
Ricci flow on certain three-manifolds. arXiv:math/0307245.

\end{thebibliography}

\end{document}
